% =====================================================================
%  CHƯƠNG 5 — KẾT LUẬN VÀ HƯỚNG PHÁT TRIỂN
% =====================================================================
\chapter{Kết luận và hướng phát triển}

\section{Kết luận}

Khóa luận đã làm được những việc cụ thể và chưa chứng minh được một điều quan trọng. Về phần đã làm:

\begin{itemize}
\item Pipeline dữ liệu SEC XBRL tái lập được: 8 công ty, 332 quý, 16 chỉ tiêu, kiểm chứng nguồn gốc
      bằng SHA-256 (0 hash lệch, 0 fact thiếu, 0 ô bịa số).
\item Bộ 47 đặc trưng tài chính gồm tỷ số thanh khoản/đòn bẩy/sinh lời/hiệu quả, tăng trưởng so với
      cùng kỳ và đặc trưng đường đi theo cửa sổ quý.
\item Bốn họ mô hình scikit-learn đặt trong một Pipeline chống rò rỉ, cộng các baseline đối chứng
      (dummy, ticker-prior, single-feature, Altman $Z''$).
\item Kết quả trên test: Random Forest được chốt đạt AUROC 0,983, AP 0,991, F1 0,960 tại ngưỡng vận
      hành 0,788.
\item Ba lớp kiểm chứng: cross-company/LOCO, kiểm định ý nghĩa (DeLong và bootstrap), và kiểm chứng
      độ nhạy theo ba định nghĩa nhãn tái lập được.
\end{itemize}

Về phần chưa chứng minh được — và đây là kết luận trung thực quan trọng nhất — trên bộ test 64 mẫu,
mô hình không vượt baseline ticker-prior về ý nghĩa thống kê, và năng lực tổng quát hoá sang công ty
chưa từng thấy còn thấp (AUROC cross-company 0,933; LOCO 0,628). Nút thắt của bài toán không nằm ở
thuật toán: đổi Random Forest sang XGBoost hay LiteSVM, hay thêm SMOTE, đều không tạo khác biệt có ý
nghĩa. Nút thắt nằm ở dữ liệu (chỉ 8 thực thể) và ở định nghĩa nhãn. Vì vậy đóng góp chính của khóa
luận được định vị là việc chỉ ra và định lượng rò rỉ thông tin ở cấp thực thể, cùng một giao thức
đánh giá đúng cho lớp bài toán này.

\section{Hướng phát triển}

\begin{enumerate}
\item \textbf{Khai thác văn bản 10-K/10-Q bằng NLP.} Phần thuyết minh MD\&A và ghi chú thường chứa
      cảnh báo định tính (going concern, kiện tụng, mất khách hàng lớn) mà các tỷ số kế toán bỏ sót.
      Đây là hướng bổ sung thông tin mạnh nhất vì nó thêm loại tín hiệu khác về bản chất, chứ không
      thêm biến thể của cùng một loại. Cách làm khả thi là mã hoá văn bản bằng mô hình ngôn ngữ, rồi
      ghép đặc trưng văn bản với 47 đặc trưng số hiện có và đánh giá lại bằng đúng khung cross-company.
\item \textbf{Thêm dữ liệu thị trường theo thời gian thực.} Giá cổ phiếu, spread tín dụng (CDS), lợi
      suất trái phiếu doanh nghiệp và biến động ngụ ý cập nhật liên tục, thay cho dữ liệu kế toán
      trễ theo quý. Kết hợp hai loại tín hiệu có thể giảm độ trễ cảnh báo — đúng vào điểm yếu số một
      đã nêu ở mục rào cản kỹ thuật.
\item \textbf{Mở rộng universe và thêm ngành.} Tăng số thực thể (nhiều công ty bán lẻ hơn, hoặc nhiều
      ngành) để cross-company học được quy luật chung thay vì nhớ mặt công ty. Đây là điều kiện cần
      để con số LOCO 0,628 được cải thiện một cách thực chất.
\item \textbf{Mô hình hoá theo thời gian (survival/sequence).} Thay phân loại từng quý bằng mô hình
      survival hoặc mô hình chuỗi khai thác tính chồng lấn của dữ liệu, dự báo thời điểm xảy ra suy
      giảm thay vì chỉ xác suất của một quý.
\item \textbf{Bổ sung sự kiện phá sản thật.} Mở rộng tập sang các công ty đã nộp 8-K item 1.03 để có
      nhãn kiểm chứng được, từ đó chuyển bài toán sang dự báo phá sản thật sự.
\end{enumerate}
